\documentclass[11pt]{article}

\usepackage[margin=1in]{geometry}
\usepackage{amsmath,amssymb}
\usepackage{booktabs}
\usepackage{graphicx}
\usepackage{xcolor}
\usepackage{hyperref}
\usepackage{enumitem}
\usepackage{caption}

\hypersetup{colorlinks=true, linkcolor=blue!50!black, urlcolor=blue!50!black}
\setlist{itemsep=2pt, topsep=4pt}
\usepackage{microtype}
\setlength{\emergencystretch}{3em}
\newcommand{\codeus}{\textunderscore\allowbreak}
\newcommand{\code}[1]{{\ttfamily\let\_\codeus #1}}

\title{Follow-up Runs of the Water-Quality Benchmark\\
       on the Colorado River Network}
\author{Andrew Chincea \\ TDA Lab}
\date{November 2026}

\begin{document}
\maketitle

\begin{abstract}
This report continues \code{kron\_vs\_full\_laplacian\_report.pdf} (the ``main report'') and collects
the benchmark runs made after it. Every run is a paired comparison against the leave-one-station-out
(LOO) benchmark of the main report (its Table~\ref{tab:base} values are repeated here), on identical
test rows. So far it covers four runs. (1) Water-year indicators on the current rows change no
target reliably except turbidity, where every linear model gains about $0.06$ and the graph model
$0.08$; pH gets worse for Elastic Net and the Spline GAM. (2) With the longer 2021--2026 history,
the water-year indicators remove the $1$\,\textdegree C cold offset found in the main report, but
the history still does not beat the current rows for any target: the records cover different years
at different stations, and for the linear models the indicators then also absorb which stations are
present, so specific conductance and pH get much worse. (3) Gradient boosting chooses the
neighbour reading for pH and turbidity only, and gains nothing reliable. (4) Hourly upstream rainfall
from NOAA MRMS radar is chosen by the inner loop in almost every fold but adds at most about $0.01$
(DO, graph and two linear models) and does not help turbidity; rain sampled at the 67 nodes of the
compressed graph does as well as rain along the whole upstream river.
\end{abstract}

% =====================================================================
\section{Setup}

The data, graph, models, nested LOO and $R^2$ (equal weight per station) are those of the main
report. Each run below changes one option of \code{water\_quality\_model\_benchmark.py} and is compared
with the same model of the main benchmark on identical test rows with \code{compare\_runs.py}:
$\Delta R^2$ is the change, the 95\% interval resamples stations with the predictions held fixed, and
better/worse counts the stations whose mean squared error fell/rose. Intervals that exclude zero are
in bold. With 17--24 stations, differences below about $0.03$--$0.05$ are within noise.

\begin{table}[h]
\centering
\small
\caption{Reference: LOO $R^2$ of the main benchmark (main report, Table~14), which every run below is
paired against.}
\label{tab:base}
\begin{tabular}{l rrrrr}
\toprule
Model & Water temp. & DO & Spec.\ cond. & pH & Turbidity \\
\midrule
Graph-Laplacian   & $0.863$ & $0.700$ & $0.385$ & $0.117$ & $0.104$ \\
OLS               & $0.863$ & $0.695$ & $-0.099$ & $-0.014$ & $0.136$ \\
Ridge             & $0.864$ & $0.701$ & $-0.110$ & $0.003$ & $0.136$ \\
Elastic Net       & $0.863$ & $0.705$ & $-0.109$ & $0.023$ & $0.136$ \\
Spline GAM        & $0.838$ & $0.708$ & $0.022$ & $0.012$ & $0.091$ \\
Gradient Boosting & $0.833$ & $0.733$ & $-0.094$ & $-0.171$ & $0.391$ \\
\bottomrule
\end{tabular}
\end{table}

% =====================================================================
\section{Water-year effect}\label{sec:wy}

\paragraph{Motivation.} In the main report (Section~15), training on the longer 2021--2026 record
made every model predict water temperature about $1$\,\textdegree C too cold: with no term for the
year, the models learned the multi-year level, while 85\% of the test rows are from 2025--26 and
water year 2026 was $1.2$\,\textdegree C warmer than usual.

\paragraph{Option.} \code{BENCHMARK\_YEAR\_EFFECT=wy} adds indicator columns for water years
2022--2026 (October--September; WY2021 is the reference, and the few hours after September 2026 count
as WY2026) to every feature candidate of every model, so the inner loop cannot drop them. Their
coefficients are estimated from the training stations, which are observed in the same years as the
held-out station, so the held-out station's values never enter. Two runs were made with the five
fast models and the graph model: on the current rows (a control, to separate the effect of the
indicators from that of the extra data) and with the longer history.

% ---------------------------------------------------------------------
\subsection{On the current rows}

\begin{table}[h]
\centering
\small
\caption{Water-year indicators on the current rows: paired $\Delta R^2$ against the main benchmark.}
\label{tab:wy}
\begin{tabular}{ll rrrcc}
\toprule
Target & Model & Without & With & $\Delta R^2$ & Interval & Better/worse \\
\midrule
Water temp.   & Graph-Laplacian & $0.863$ & $0.859$ & $-0.005$ & ($-0.03$, $0.01$) & 12/12 \\
              & OLS             & $0.863$ & $0.862$ & $-0.001$ & ($-0.02$, $0.01$) & 14/10 \\
              & Ridge           & $0.864$ & $0.863$ & $-0.001$ & ($-0.02$, $0.01$) & 13/11 \\
              & Elastic Net     & $0.863$ & $0.866$ & $+0.002$ & ($-0.02$, $0.02$) & 15/9 \\
              & Spline GAM      & $0.838$ & $0.820$ & $\mathbf{-0.018}$ & ($-0.03$, $-0.00$) & 7/17 \\
\addlinespace
DO            & Graph-Laplacian & $0.700$ & $0.694$ & $-0.005$ & ($-0.03$, $0.02$) & 6/11 \\
              & OLS             & $0.695$ & $0.689$ & $-0.006$ & ($-0.03$, $0.02$) & 9/8 \\
              & Ridge           & $0.701$ & $0.701$ & $+0.000$ & ($-0.02$, $0.03$) & 8/9 \\
              & Elastic Net     & $0.705$ & $0.702$ & $-0.003$ & ($-0.02$, $0.01$) & 8/9 \\
              & Spline GAM      & $0.708$ & $0.695$ & $\mathbf{-0.013}$ & ($-0.03$, $-0.00$) & 5/12 \\
\addlinespace
Spec.\ cond.  & Graph-Laplacian & $0.385$ & $0.393$ & $+0.009$ & ($-0.08$, $0.15$) & 12/12 \\
              & OLS             & $-0.099$ & $-0.025$ & $+0.074$ & ($-0.14$, $0.30$) & 10/14 \\
              & Ridge           & $-0.110$ & $-0.073$ & $+0.037$ & ($-0.13$, $0.16$) & 11/13 \\
              & Elastic Net     & $-0.109$ & $-0.070$ & $+0.039$ & ($-0.12$, $0.16$) & 11/13 \\
              & Spline GAM      & $0.022$ & $-0.037$ & $-0.059$ & ($-0.26$, $0.03$) & 9/15 \\
\addlinespace
pH            & Graph-Laplacian & $0.117$ & $0.132$ & $+0.015$ & ($-0.00$, $0.04$) & 10/9 \\
              & OLS             & $-0.014$ & $-0.054$ & $-0.040$ & ($-0.09$, $0.01$) & 8/11 \\
              & Ridge           & $0.003$ & $-0.003$ & $-0.006$ & ($-0.04$, $0.04$) & 8/11 \\
              & Elastic Net     & $0.023$ & $-0.058$ & $\mathbf{-0.081}$ & ($-0.12$, $-0.05$) & 4/15 \\
              & Spline GAM      & $0.012$ & $-0.066$ & $\mathbf{-0.078}$ & ($-0.15$, $-0.02$) & 6/13 \\
\addlinespace
Turbidity     & Graph-Laplacian & $0.104$ & $0.186$ & $+0.081$ & ($-0.04$, $0.47$) & 12/5 \\
              & OLS             & $0.136$ & $0.194$ & $\mathbf{+0.058}$ & ($0.01$, $0.22$) & 13/4 \\
              & Ridge           & $0.136$ & $0.196$ & $\mathbf{+0.060}$ & ($0.01$, $0.22$) & 13/4 \\
              & Elastic Net     & $0.136$ & $0.193$ & $\mathbf{+0.056}$ & ($0.01$, $0.20$) & 13/4 \\
              & Spline GAM      & $0.091$ & $0.135$ & $+0.044$ & ($-0.06$, $0.14$) & 12/5 \\
\bottomrule
\end{tabular}
\end{table}

\begin{itemize}
  \item \textbf{Water temperature, DO and specific conductance do not change} beyond noise for the
        graph and linear models. On the current rows the training and test stations cover the same
        months, so there is little common offset to remove (mean signed station bias $-0.05$ before,
        $+0.05$ to $+0.09$\,\textdegree C after). The Spline GAM loses a little on water temperature
        and DO.
  \item \textbf{Turbidity improves for every model}: the three linear models gain about $0.06$ with
        intervals above zero (13 of 17 stations better), and the graph model gains $0.08$
        ($0.104\to0.186$, interval crossing zero). Turbidity depends on runoff, which differs a lot
        from year to year, so a year level shared across stations is plausible. The best turbidity
        model is still gradient boosting ($0.391$), which was not part of this run.
  \item \textbf{pH gets worse} for Elastic Net and the Spline GAM ($-0.08$, intervals below zero)
        and does not change for the graph model ($+0.015$).
  \item The indicators were fixed in advance (water years, no tuning), so these are fair results.
        Because they help only turbidity, the option stays off by default.
\end{itemize}

% ---------------------------------------------------------------------
\subsection{With the longer history}

\begin{table}[h]
\centering
\small
\caption{Water-year indicators with the 2021--2026 history: LOO $R^2$ of the main benchmark (Base),
the history run of the main report (History) and the history run with water-year indicators
(History + WY). $\Delta$: History + WY against History, with the paired 95\% interval.}
\label{tab:histwy}
\begin{tabular}{ll rrr rc}
\toprule
Target & Model & Base & History & History + WY & $\Delta$ & Interval \\
\midrule
Water temp.   & Graph-Laplacian & $0.863$ & $0.816$ & $0.813$ & $-0.002$ & ($-0.04$, $0.03$) \\
              & OLS             & $0.863$ & $0.840$ & $0.856$ & $+0.016$ & ($-0.01$, $0.04$) \\
              & Ridge           & $0.864$ & $0.840$ & $0.850$ & $+0.009$ & ($-0.02$, $0.03$) \\
              & Elastic Net     & $0.863$ & $0.841$ & $0.851$ & $+0.010$ & ($-0.01$, $0.03$) \\
              & Spline GAM      & $0.838$ & $0.745$ & $0.757$ & $+0.012$ & ($-0.05$, $0.08$) \\
\addlinespace
DO            & Graph-Laplacian & $0.700$ & $0.699$ & $0.706$ & $+0.007$ & ($-0.02$, $0.04$) \\
              & OLS             & $0.695$ & $0.702$ & $0.694$ & $-0.008$ & ($-0.03$, $0.01$) \\
              & Ridge           & $0.701$ & $0.708$ & $0.698$ & $-0.010$ & ($-0.03$, $0.01$) \\
              & Elastic Net     & $0.705$ & $0.712$ & $0.708$ & $-0.004$ & ($-0.02$, $0.01$) \\
              & Spline GAM      & $0.708$ & $0.719$ & $0.715$ & $-0.004$ & ($-0.03$, $0.01$) \\
\addlinespace
Spec.\ cond.  & Graph-Laplacian & $0.385$ & $0.334$ & $0.322$ & $-0.013$ & ($-0.15$, $0.04$) \\
              & OLS             & $-0.099$ & $-0.022$ & $-0.727$ & $-0.706$ & ($-4.08$, $0.08$) \\
              & Ridge           & $-0.110$ & $-0.032$ & $-0.227$ & $-0.195$ & ($-1.29$, $0.08$) \\
              & Elastic Net     & $-0.109$ & $-0.039$ & $-0.124$ & $-0.085$ & ($-0.67$, $0.07$) \\
              & Spline GAM      & $0.022$ & $0.055$ & $0.066$ & $+0.011$ & ($-0.12$, $0.06$) \\
\addlinespace
pH            & Graph-Laplacian & $0.117$ & $-0.037$ & $-0.051$ & $-0.013$ & ($-0.08$, $0.04$) \\
              & OLS             & $-0.014$ & $-0.011$ & $-0.273$ & $\mathbf{-0.262}$ & ($-0.66$, $-0.03$) \\
              & Ridge           & $0.003$ & $0.013$ & $-0.173$ & $\mathbf{-0.186}$ & ($-0.51$, $-0.02$) \\
              & Elastic Net     & $0.023$ & $0.029$ & $0.006$ & $-0.023$ & ($-0.08$, $0.00$) \\
              & Spline GAM      & $0.012$ & $0.009$ & $0.006$ & $-0.003$ & ($-0.03$, $0.03$) \\
\addlinespace
Turbidity     & Graph-Laplacian & $0.104$ & $0.043$ & $0.095$ & $\mathbf{+0.052}$ & ($0.02$, $0.13$) \\
              & OLS             & $0.136$ & $0.032$ & $0.100$ & $\mathbf{+0.068}$ & ($0.01$, $0.12$) \\
              & Ridge           & $0.136$ & $0.031$ & $0.095$ & $+0.064$ & ($-0.00$, $0.10$) \\
              & Elastic Net     & $0.136$ & $0.033$ & $0.099$ & $+0.067$ & ($-0.01$, $0.10$) \\
              & Spline GAM      & $0.091$ & $0.085$ & $0.128$ & $+0.043$ & ($-0.01$, $0.09$) \\
\bottomrule
\end{tabular}
\end{table}

\begin{table}[h]
\centering
\small
\caption{Water temperature: decomposition of the station-level error. Signed bias: mean over
stations of the station-mean error (\textdegree C; without 09429490). Bias$^2$: mean squared
station-mean error. Within: mean variance of the error about its station mean. $R^2$ without
09429490 in the last column.}
\label{tab:wtdecomp}
\begin{tabular}{ll rrrr}
\toprule
Model & Run & Signed bias & Bias$^2$ & Within & $R^2$ excl. \\
\midrule
Graph-Laplacian & Base         & $-0.05$ & $0.90$ & $5.96$ & $0.855$ \\
                & History      & $-0.92$ & $3.25$ & $6.01$ & $0.833$ \\
                & History + WY & $+0.04$ & $3.09$ & $6.28$ & $0.840$ \\
\addlinespace
Ridge           & Base         & $-0.05$ & $0.88$ & $5.95$ & $0.856$ \\
                & History      & $-0.98$ & $2.01$ & $6.02$ & $0.831$ \\
                & History + WY & $+0.01$ & $1.27$ & $6.29$ & $0.841$ \\
\bottomrule
\end{tabular}
\end{table}

\begin{itemize}
  \item \textbf{The cold offset is gone} (Table~\ref{tab:wtdecomp}): the mean signed station bias
        of water temperature goes from about $-0.95$ to $+0.01$--$+0.04$\,\textdegree C, and the
        linear models gain about $0.01$ ($0.840\to0.850$--$0.856$).
  \item \textbf{But the history still does not beat the current rows.} Water temperature stays
        $0.01$--$0.05$ below the main benchmark for every model, and the graph model does not
        change at all ($0.816\to0.813$). Two things remain after the offset is removed: the station
        biases are still more spread out than in the main benchmark (bias$^2$ $1.27$ against
        $0.88$ for Ridge; for the graph model most of the $3.09$ is the 09429490 fold, whose bias is
        still about $-6$\,\textdegree C), and the within-station error is larger ($5.96\to6.28$).
        The seasonal coefficients fitted over six years describe the 2025--26 test period less
        well than those fitted on the same period.
  \item \textbf{Specific conductance and pH get much worse for OLS and Ridge.} OLS falls to $-0.73$
        for specific conductance and $-0.27$ for pH, Ridge to $-0.23$ and $-0.17$ (pH intervals
        below zero). The specific-conductance loss comes from a few stations, mostly 09510000 (Verde
        below Bartlett, predicted $1{,}100$--$1{,}900$\,\textmu S/cm too high) and the
        high-conductance station 96751415 (about $1{,}150$\,\textmu S/cm too low). The graph model and the Spline GAM barely change.
  \item \textbf{Likely cause: the records cover different years at different stations.} In the
        history tables, some stations report only in the last one or two water years (09510000 from
        WY2024, 92116891 from WY2025, 75465016 only in WY2026) and others only in the first ones
        (133114453 and 134075067 stop in WY2022 for specific conductance;
        \code{results/history\_coverage\_by\_wy.csv}). The linear models have no term for each
        station, so a water-year coefficient also absorbs which stations report in that year, and
        that mix changes with every held-out station. For specific conductance, whose level differs
        by a factor of five between stations, this swamps the year effect. The graph model has a
        level $g_i$ for every station, so its year coefficients are estimated within stations,
        which is why it is unaffected. This explanation was not tested separately.
  \item \textbf{Turbidity improves against the history run} ($+0.04$ to $+0.07$, intervals above
        zero for the graph model and OLS), consistent with the current-row run, but stays at or
        below the main benchmark.
  \item \textbf{Conclusion.} With the year effect, the longer history no longer biases water
        temperature, but it still does not improve any target over the main benchmark, and with
        the uneven coverage it makes the linear models unstable. A fair use of the history would
        need the year effect estimated within stations (station-level offsets for every model, as
        the graph model has), or training restricted to stations that report in every year. Neither
        is run yet. The history option stays off by default.
\end{itemize}

% =====================================================================
\section{Gradient boosting with the neighbour reading}\label{sec:nbrgb}

The neighbour reading of the main report (Section~14: the mean same-hour anomaly of the target at
the three nearest training stations) was run for gradient boosting, which was missing from that
section. As before, every feature candidate is offered with and without it and the inner loop
chooses.

\begin{table}[h]
\centering
\small
\caption{Gradient boosting with the neighbour reading: paired $\Delta R^2$ against gradient boosting
in the main benchmark. Folds: outer folds in which the inner loop chose a candidate with the reading.}
\label{tab:nbrgb}
\begin{tabular}{l rrrcc c}
\toprule
Target & Without & With & $\Delta R^2$ & Interval & Better/worse & Folds \\
\midrule
Water temp.  & $0.833$ & $0.833$ & $0.000$ & -- & 0/0 & 0/24 \\
DO           & $0.733$ & $0.729$ & $-0.004$ & ($-0.01$, $0.00$) & 0/1 & 1/17 \\
Spec.\ cond. & $-0.094$ & $-0.094$ & $0.000$ & -- & 0/0 & 0/24 \\
pH           & $-0.171$ & $-0.123$ & $+0.048$ & ($-0.04$, $0.17$) & 11/6 & 17/19 \\
Turbidity    & $0.391$ & $0.375$ & $-0.016$ & ($-0.10$, $0.01$) & 8/8 & 16/17 \\
\bottomrule
\end{tabular}
\end{table}

\begin{itemize}
  \item \textbf{Gradient boosting uses the reading only for pH and turbidity}, and never for water
        temperature or specific conductance, unlike the linear models, which chose it in every fold
        for every target but turbidity.
  \item \textbf{pH improves by $0.05$, but stays below zero} ($-0.123$) and the interval crosses
        zero. The linear models gained $0.12$--$0.14$ from the same reading and the graph model
        reached $0.178$, so the reading does not make gradient boosting competitive for pH.
  \item \textbf{Turbidity is unchanged} ($-0.016$, within noise), so gradient boosting without the
        reading ($0.391$) remains the best turbidity model.
\end{itemize}

% =====================================================================
\section{Upstream rainfall}\label{sec:rain}

\paragraph{Data.} Hourly NOAA MRMS radar--gauge precipitation (MultiSensor QPE, 1-hour, Pass~2;
public on AWS, $0.01^\circ$ grid, from October 2020), downloaded by \code{download\_mrms.py} for
every station's benchmark window at the $7{,}259$ grid cells that contain a node of the uncompressed
river graph. 83 hours are missing from the archive, none in the main 2025--26 window. Radar
undercounts snow in the mountains, so winter precipitation in the headwaters is low.

\paragraph{Features.} For each station and hour, \code{build\_rain\_features.py} averages the rain
over the river upstream of the station (its ancestors on the directed graph, plus the station)
and sums it over the last $1$, $6$, $24$, $72$ and $168$ hours, as $\log(1+\text{mm})$. Every
window ends before the hour of the reading, so no later rain is used. The windows were fixed in
advance. Two ways of sampling the upstream rain were compared:
\begin{description}[leftmargin=1.5em]
  \item[uncompressed] every node of the 90\,m river graph (median 187 grid cells per station), so the
        mean weights each cell by the length of upstream channel in it;
  \item[compressed] only the 67 nodes of the compressed graph (stations, gauges, confluences and
        channel ends; median 6 cells per station).
\end{description}
The two schemes correlate at $0.90$--$0.97$. Within a station, \code{rain\_72h} correlates $0.35$
with turbidity and about $0.01$ with specific conductance.

\paragraph{Option.} \code{BENCHMARK\_RAIN=uncompressed|compressed} offers every feature candidate
with and without the five rain columns, and the inner loop chooses, as for the neighbour reading.
Both runs used the five fast models and the graph model; gradient boosting and random forest were
not part of these runs.

\begin{table}[h]
\centering
\small
\caption{Upstream rainfall: paired $\Delta R^2$ against the main benchmark, for rain sampled on the
uncompressed and on the compressed graph. Folds: outer folds whose chosen candidate included rain
(uncompressed/compressed).}
\label{tab:rain}
\footnotesize
\setlength{\tabcolsep}{4pt}
\begin{tabular}{ll r rc rc c}
\toprule
 & & & \multicolumn{2}{c}{Uncompressed} & \multicolumn{2}{c}{Compressed} & \\
\cmidrule(lr){4-5}\cmidrule(lr){6-7}
Target & Model & Base & $\Delta R^2$ & Interval & $\Delta R^2$ & Interval & Folds \\
\midrule
Water temp.   & Graph-Laplacian & $0.863$ & $0.000$ & ($-0.00$, $0.00$) & $+0.001$ & ($-0.00$, $0.00$) & 24/24 of 24 \\
              & OLS             & $0.863$ & $+0.001$ & ($-0.00$, $0.00$) & $+0.001$ & ($-0.00$, $0.00$) & 24/24 of 24 \\
              & Ridge           & $0.864$ & $+0.001$ & ($-0.00$, $0.00$) & $+0.001$ & ($-0.00$, $0.00$) & 24/23 of 24 \\
              & Elastic Net     & $0.863$ & $0.000$ & ($-0.00$, $0.00$) & $0.000$ & ($-0.00$, $0.00$) & 21/20 of 24 \\
              & Spline GAM      & $0.838$ & $+0.005$ & ($-0.00$, $0.02$) & $+0.005$ & ($-0.00$, $0.02$) & 24/24 of 24 \\
\addlinespace
DO            & Graph-Laplacian & $0.700$ & $\mathbf{+0.012}$ & ($0.00$, $0.02$) & $\mathbf{+0.014}$ & ($0.01$, $0.03$) & 17/17 of 17 \\
              & OLS             & $0.695$ & $+0.004$ & ($-0.01$, $0.01$) & $+0.007$ & ($-0.00$, $0.01$) & 17/17 of 17 \\
              & Ridge           & $0.701$ & $+0.001$ & ($-0.01$, $0.01$) & $\mathbf{+0.011}$ & ($0.00$, $0.02$) & 17/17 of 17 \\
              & Elastic Net     & $0.705$ & $0.000$ & ($-0.01$, $0.01$) & $\mathbf{+0.007}$ & ($0.00$, $0.01$) & 17/17 of 17 \\
              & Spline GAM      & $0.708$ & $-0.004$ & ($-0.01$, $0.00$) & $+0.001$ & ($-0.00$, $0.01$) & 8/17 of 17 \\
\addlinespace
Spec.\ cond.  & Graph-Laplacian & $0.385$ & $+0.003$ & ($-0.00$, $0.01$) & $+0.002$ & ($-0.00$, $0.01$) & 24/24 of 24 \\
              & OLS             & $-0.099$ & $+0.022$ & ($-0.03$, $0.11$) & $+0.014$ & ($-0.02$, $0.09$) & 24/24 of 24 \\
              & Ridge           & $-0.110$ & $+0.012$ & ($-0.05$, $0.09$) & $-0.004$ & ($-0.05$, $0.06$) & 24/24 of 24 \\
              & Elastic Net     & $-0.109$ & $+0.013$ & ($-0.04$, $0.09$) & $+0.006$ & ($-0.02$, $0.07$) & 24/24 of 24 \\
              & Spline GAM      & $0.022$ & $-0.010$ & ($-0.06$, $0.01$) & $-0.006$ & ($-0.03$, $0.01$) & 19/4 of 24 \\
\addlinespace
pH            & Graph-Laplacian & $0.117$ & $-0.033$ & ($-0.16$, $0.02$) & $-0.030$ & ($-0.16$, $0.02$) & 19/19 of 19 \\
              & OLS             & $-0.014$ & $+0.003$ & ($-0.03$, $0.03$) & $+0.005$ & ($-0.02$, $0.03$) & 19/19 of 19 \\
              & Ridge           & $0.003$ & $+0.012$ & ($-0.01$, $0.03$) & $+0.012$ & ($-0.01$, $0.03$) & 19/19 of 19 \\
              & Elastic Net     & $0.023$ & $0.000$ & -- & $0.000$ & -- & 0/0 of 19 \\
              & Spline GAM      & $0.012$ & $+0.015$ & ($-0.03$, $0.06$) & $+0.011$ & ($-0.03$, $0.05$) & 18/18 of 19 \\
\addlinespace
Turbidity     & Graph-Laplacian & $0.104$ & $-0.011$ & ($-0.10$, $0.02$) & $-0.011$ & ($-0.10$, $0.02$) & 15/16 of 17 \\
              & OLS             & $0.136$ & $0.000$ & -- & $-0.002$ & ($-0.02$, $0.00$) & 0/1 of 17 \\
              & Ridge           & $0.136$ & $0.000$ & -- & $-0.002$ & ($-0.02$, $0.00$) & 0/1 of 17 \\
              & Elastic Net     & $0.136$ & $0.000$ & -- & $-0.002$ & ($-0.02$, $0.00$) & 0/1 of 17 \\
              & Spline GAM      & $0.091$ & $0.000$ & ($-0.02$, $0.02$) & $-0.004$ & ($-0.04$, $0.01$) & 15/15 of 17 \\
\bottomrule
\end{tabular}
\end{table}

\begin{itemize}
  \item \textbf{The inner loop chooses rain almost everywhere, but it changes little.} For water
        temperature, DO, specific conductance and pH, nearly every outer fold of every model picks a
        candidate with rain, yet all changes are below $0.025$. The only intervals above zero are for
        DO: the graph model gains $0.012$--$0.014$ with both schemes (11--12 of 17 stations better),
        and Ridge and Elastic Net gain $0.011$ and $0.007$ with the compressed scheme. These are
        small.
  \item \textbf{Turbidity does not improve}, although it is the target most related to rain within
        a station. The linear models never choose rain under the uncompressed scheme (once under
        the compressed one). The graph model and the Spline GAM do choose it, but the graph model
        loses $0.01$ and the Spline GAM gains nothing. Turbidity $R^2$ is dominated by a few stations with very large and
        variable readings: the graph model's loss comes mostly from 09365000 (San Juan River at
        Farmington), whose mean bias grows from $618$ to $691$\,FNU. A possible reason is that sediment jumps after storms above some threshold, which a linear
        term in log rain may not capture;
        gradient boosting, the best turbidity model ($0.391$), was not run with rain.
  \item \textbf{The graph model's pH loss ($-0.03$) comes from one station}, 09429490 above
        Imperial Dam, the outlying fold of the main report, whose mean bias moves from $-0.08$ to
        $-0.23$. Of the other 18 stations, 15 (uncompressed) and 16 (compressed) are better; without
        09429490 the graph model's pH $R^2$ rises from $0.123$ to $0.135$ (uncompressed) and $0.139$
        (compressed).
  \item \textbf{Uncompressed and compressed rain give the same results.} Paired directly, no
        difference exceeds $0.016$ for any model or target. The only intervals above zero favour the
        compressed scheme for DO (Elastic Net $+0.007$, Spline GAM $+0.005$; Ridge $+0.010$ and the
        graph model $+0.002$ with intervals touching zero). Rain at the 67 nodes of the compressed graph
        carries as much information as rain along every 90\,m cell upstream.
\end{itemize}

% =====================================================================
\section{Conclusions so far}

\begin{enumerate}
  \item \textbf{A water-year effect helps turbidity only.} On the current rows it adds about $0.06$
        to every linear model (intervals above zero) and $0.08$ to the graph model, changes nothing
        for water temperature, DO and specific conductance, and lowers pH for Elastic Net and the
        Spline GAM.
  \item \textbf{The longer history is still not useful.} The year effect removes its
        $1$\,\textdegree C water-temperature offset, but no target beats the main benchmark, and
        because the stations' records cover different years the linear models become unstable for
        specific conductance and pH. The graph model, which has a level for every station, is
        unaffected by that problem but gains nothing either.
  \item \textbf{The neighbour reading does not help gradient boosting} for any target beyond noise.
        Gradient boosting stays the best turbidity model without it.
  \item \textbf{Upstream rainfall adds little.} It is chosen in nearly every fold, but the only
        reliable gains are about $0.01$ for DO; turbidity, the target that rain should help most, does
        not improve for the graph or linear models. Sampling rain at the compressed graph's 67 nodes
        is as good as sampling it along every upstream cell.
  \item The main benchmark (Table~\ref{tab:base}) remains the reference.
\end{enumerate}

\paragraph{Next steps (proposed, not run).} Turbidity with water-year indicators or upstream rain
for gradient boosting and the Graph + GB hybrid (rain enters turbidity through thresholds that a
tree model can represent); the history with station-level offsets for the linear models or
restricted to stations that report in every year; clipping turbidity at the sensor maximum
(3{,}800--4{,}000\,FNU in the added years); a dam-tailwater indicator for the Lower Colorado.

\paragraph{Reproducibility.} The runs come from \code{water\_quality\_model\_benchmark.py} at commit
\code{9f26e65} on the lab server, with
\code{BENCHMARK\_MODELS=Dummy,OLS,Ridge,Elastic Net,Spline GAM,Graph-Laplacian} and
\code{BENCHMARK\_YEAR\_EFFECT=wy}, without and with \code{BENCHMARK\_HISTORY=data/benchmark\_data\_history},
saved in \code{results/benchmark\_results\_wy/} and \code{results/benchmark\_results\_history\_wy/}, and
with \code{BENCHMARK\_MODELS=Gradient Boosting} and \code{BENCHMARK\_NEIGHBOUR=on}, saved in
\code{results/benchmark\_results\_neighbour\_gb/}. The server logs are \code{wy.log},
\code{hist\_wy.log} and \code{nbr\_gb.log}. The paired comparisons come from \code{compare\_runs.py} against
\code{results/benchmark\_results\_inner\_fs\_graph/} (graph model),
\code{results/benchmark\_results\_inner\_fs/} (other models) and the history run of the main report.
They are saved as \code{results/wy\_vs\_base\_*.csv}, \code{results/history\_wy\_vs\_base\_*.csv},
\code{results/history\_wy\_vs\_history\_*.csv} and
\code{results/neighbour\_vs\_base\_gradient\_boosting.csv}.
The rainfall runs use \code{BENCHMARK\_RAIN=uncompressed|compressed} at commit \code{b2230e5}
(one process per target, \code{run\_rain\_parallel.sh}, merged by \code{merge\_target\_parts.py}),
with the same six models, saved in \code{results/benchmark\_results\_rain\_uncompressed/} and
\code{results/benchmark\_results\_rain\_compressed/}; the rain comes from \code{download\_mrms.py}
(\code{data/rain/mrms/}) and \code{build\_rain\_features.py}
(\code{data/benchmark\_data/rain\_features\_*.pkl}). Their comparisons are
\code{results/rain\_uncompressed\_vs\_base\_*.csv}, \code{results/rain\_compressed\_vs\_base\_*.csv} and
\code{results/rain\_compressed\_vs\_uncompressed\_*.csv}.
The error decomposition and the coverage table come from \code{year\_effect\_analysis.py} (seconds),
saved in \code{results/year\_effect\_diagnostics.csv} and \code{results/history\_coverage\_by\_wy.csv}.

\end{document}
